\chapter*{Glossario}
\markboth{Glossario}{Glossario}
\addcontentsline{toc}{chapter}{Glossario}

Questo glossario raccoglie i termini tecnici usati nel libro. Le epoche preistoriche riportano sempre le date nella convenzione adottata nel testo: anni prima dell'anno zero astronomico, ottenute sottraendo circa 1.950 anni dai valori cal BP originali della letteratura.

\begin{description}

\item[Anni cal BP] Anni trascorsi rispetto al 1950, ricavati da datazioni al radiocarbonio e corretti (calibrati) per le variazioni storiche del tenore di $^{14}$C nell'atmosfera. La calibrazione è necessaria perché la produzione di $^{14}$C non è costante nel tempo: senza correzione, le date radiocarboniche risulterebbero sistematicamente più giovani del reale. La sigla BP sta per \textit{Before Present}, dove il ``presente'' è convenzionalmente fissato al 1950. In questo libro le date sono sempre espresse in anni prima dell'anno zero astronomico, sottraendo circa 1.950 anni dal valore cal BP.

\item[ARC model] Modello del costo energetico della locomozione terrestre in funzione della pendenza, proposto da Herman Pontzer (\textit{Biology Letters}, 2016). Il nome deriva dai due processi della fisiologia muscolare su cui si fonda: \textit{Activation-Relaxation} (attivazione e rilassamento delle fibre muscolari) e \textit{Cross-bridge cycling} (ciclo dei ponti actina-miosina). Il modello è ricavato dai meccanismi cellulari della contrazione, non da una curva adattata empiricamente, ed è stato validato su un intervallo di masse corporee da 0,78 grammi a 431 chilogrammi e su pendenze dalla discesa moderata all'arrampicata verticale, con buon accordo tra le previsioni e i dati misurati per specie molto diverse. Per le specie di interesse di questo libro (camoscio, stambecco, cervo, essere umano), tutte con massa corporea tra quaranta e novanta chilogrammi, il modello si trova nel centro del suo intervallo di validazione.

\item[Bølling-Allerød] Fase di miglioramento climatico rapido compresa tra circa 12.700 e 11.000 anni prima dell'anno zero, intercalata all'interno del Tardoglaciale. Le temperature medie estive sull'arco alpino aumentarono di diversi gradi in poche centinaia di anni; le foreste avanzarono rapidamente in quota e il limite del bosco si spostò verso l'alto di centinaia di metri. Il nome è composto da quello di due località del Nord Europa — Bølling in Danimarca e Allerød anch'essa danese — dove le sequenze polliniche che hanno permesso di identificare questa fase furono descritte per la prima volta nella prima metà del Novecento. La fase si conclude con il brusco raffreddamento del Dryas Recente.

\item[Cal BP] Vedi \textbf{Anni cal BP}.

\item[Castelnoviano] Cultura litica che contraddistingue la fase più recente del Mesolitico nell'Italia settentrionale, collocata approssimativamente tra 9.500 e 7.500 anni prima dell'anno zero. Prende il nome dal sito di Castellnovo d'Elsa (Siena), dove il materiale litico rappresentativo di questa fase fu descritto per la prima volta. Si distingue dal Sauveterriano precedente per la presenza di trapezi e lame più regolari nell'industria litica, e per criteri insediativi in parte diversi: i siti castelnoviani mostrano una maggiore preferenza per l'esposizione solare rispetto ai soli criteri di visibilità e controllo del territorio.

\item[Crioclastismo] Processo fisico di frattura della roccia causato dalla ripetuta alternanza di gelo e disgelo dell'acqua infiltrata nelle fessure. L'acqua, ghiacciando, aumenta di volume di circa il 9\% ed esercita pressioni che possono superare la resistenza alla trazione della roccia, provocandone la frantumazione e il distacco di frammenti. Il crioclastismo è particolarmente intenso durante i periodi di transizione climatica, quando il numero di cicli gelo-disgelo è più elevato. Alla Grotta di Fumane, in Lessinia, il distacco dei frammenti di parete recanti pitture paleolitiche è attribuito dagli scavatori a questo meccanismo.

\item[Criterio di informazione (AIC, BIC)] Misure statistiche per confrontare modelli concorrenti costruiti sugli stessi dati, che bilanciano la bontà di adattamento con la complessità del modello. Il criterio di Akaike (AIC, Hirotugu Akaike, 1974) penalizza ogni parametro aggiuntivo con un costo fisso pari a 2; il criterio bayesiano (BIC, Gideon Schwarz, 1978) applica una penalità che cresce con il logaritmo della dimensione del campione, risultando più severo con dataset grandi. Un modello con AIC o BIC più basso è preferibile a parità di tutte le altre condizioni. Entrambi i criteri sono usati per confrontare modelli spaziali nella letteratura sui siti mesolitici alpini.

\item[Dryas Recente] Fase di raffreddamento climatico brusca e marcata, compresa tra circa 11.000 e 9.700 anni prima dell'anno zero, che interruppe il miglioramento del Bølling-Allerød. Sull'arco alpino le temperature estive calarono di diversi gradi; in alcune zone i ghiacciai avanzarono nuovamente. La durata, circa 1.300 anni, la distingue dai raffreddamenti minori del Tardoglaciale. Il nome deriva da \textit{Dryas octopetala}, una pianta erbacea artica i cui pollini sono abbondanti nelle sequenze sedimentarie di questa fase e servono da indicatore paleoclimatico. Si conclude con il passaggio rapido all'Olocene, circa 9.700 anni prima dell'anno zero.

\item[Epigravettiano] Cultura litica diffusa in Italia e nelle aree adriatiche dell'Europa meridionale nella fase finale del Paleolitico superiore, approssimativamente tra 23.000 e 11.000 anni prima dell'anno zero. Il nome allude a una derivazione dal Gravettiano, cultura più antica. L'industria è caratterizzata da lamelle a dorso, punte a dorso e geometrici; nella fase più recente compaiono elementi che anticipano le industrie mesolitiche. Ai siti epigravettiani veneti appartengono, tra gli altri, Riparo Tagliente (fase superiore), Villabruna e il riparo Dalmeri, quest'ultimo datato con precisione al periodo del Bølling-Allerød.

\item[Home range] Area geografica abitualmente utilizzata da un individuo o da un gruppo per svolgere le proprie attività vitali (alimentazione, riproduzione, riposo), nell'arco di un anno o di una stagione. Non include spostamenti eccezionali né dispersioni. A differenza del territorio in senso etologico, l'home range non implica difesa attiva dai conspecifici. Per le specie alpine discusse in questo libro i valori tipici, misurati con radiotelemetria e GPS, sono: camoscio 200--600 ettari (stagionale), stambecco 300--800 ettari, cervo 1.000--5.000 ettari con variabilità legata al sesso e alla stagione.

\item[Laacher See Tephra] Strato di cenere vulcanica emessa dall'eruzione del vulcano Laacher See (odierna Renania-Palatinato, Germania) circa 11.050--10.950 anni prima dell'anno zero. L'eruzione, classificata come pliniana con indice di esplosività vulcanica 6, disperse materiale piroclastico su un'area di oltre un milione di chilometri quadrati. Lo strato di tephra, sottile e chimicamente omogeneo, si riconosce nelle sequenze sedimentarie lacustri e torbose dell'Europa centrale e settentrionale, incluse alcune località del nord Italia. Poiché si è depositato in un intervallo di anni o decenni su un'area vastissima, costituisce un livello isocronico usato dai geologi per sincronizzare stratigrafie di regioni diverse.

\item[Megaloceros giganteus] Cervide estinto, noto comunemente come cervo gigante o alce irlandese. Caratterizzato da palchi di dimensioni eccezionali, che potevano raggiungere i 3,7 metri di apertura, su un corpo di taglia paragonabile a quella di un alce moderno. Era diffuso dall'Irlanda all'Asia centrale durante il Pleistocene. Le ultime popolazioni si estinsero circa 7.700 anni prima dell'anno zero. Nei depositi del Paleolitico medio e superiore dei siti prealpini veneti i suoi resti compaiono sporadicamente, in misura nettamente inferiore rispetto ai cervidi di taglia più piccola.

\item[Mesolitico] Periodo culturale compreso tra la fine dell'ultima glaciazione e la comparsa dell'agricoltura in una data regione. Nelle Alpi orientali si colloca approssimativamente tra 12.700 e 7.500 anni prima dell'anno zero, corrispondendo al Tardoglaciale finale e all'Olocene antico. È caratterizzato da un'economia di caccia e raccolta con industria litica microlitica (armature geometriche, trapezi, microliti) e da una progressiva specializzazione nella caccia agli ungulati di montagna, in particolare stambecco e camoscio. Si suddivide convenzionalmente in Sauveterriano (fase più antica) e Castelnoviano (fase più recente).

\item[Modelli a pattern di punti (\textit{point-process models})] Famiglia di modelli statistici per l'analisi della distribuzione spaziale di eventi puntiformi (presenze di specie, siti archeologici, episodi di incendio) in un'area continua. Un modello di processo puntuale specifica come l'intensità attesa degli eventi varia nello spazio in funzione di covariate ambientali o di altra natura. Tra i modelli più usati in archeologia e paleoecologia figurano i processi di Poisson non omogenei e le loro estensioni. Permettono di confrontare, mediante criteri come AIC e BIC, l'importanza relativa di fattori concorrenti (morfologia, clima, sforzo di ricerca) nel determinare la distribuzione osservata dei siti noti.

\item[Neanderthal (\textit{Homo neanderthalensis})] Specie del genere \textit{Homo}, evolutasi in Europa e Asia occidentale a partire da circa 400.000 anni fa. Anatomicamente distinta da \textit{Homo sapiens} per il cranio basso e allungato, le arcate sopraorbitali prominenti, il prognatismo medio-facciale e la robustezza dello scheletro postcraniale. Capace di comportamenti culturalmente complessi: produzione di strumenti musteriani, uso del fuoco, sepolture intenzionali, e in alcuni contesti uso di pigmenti e ornamenti. Si estinse circa 38.000 anni prima dell'anno zero, in un periodo di sovrapposizione con \textit{Homo sapiens} nella penisola iberica e in Italia. I siti prealpini veneti attribuiti a Neanderthal --- tra cui il Riparo Tagliente nella sua fase più antica e i siti della Lessinia --- risalgono a un periodo compreso tra 58.000 e 38.000 anni prima dell'anno zero.

\item[Neolitico] Periodo culturale caratterizzato dalla comparsa e dalla diffusione dell'economia produttiva (agricoltura e allevamento), della ceramica e, in molte regioni, dell'insediamento stabile. In Europa centro-meridionale si diffonde a partire da circa 9.000--8.000 anni prima dell'anno zero, con un avanzamento dalle aree balcaniche e danubiane verso ovest e nord. In Veneto le prime evidenze agricole sono datate intorno a 8.000 anni prima dell'anno zero. Il Neolitico segna la fine del Mesolitico e l'inizio di una trasformazione progressiva del paesaggio per intervento umano diretto, inclusi il disboscamento, la pastorizia d'altura e le prime forme di arte rupestre in contesto alpino non più strettamente legato alla caccia.

\item[Paleolitico] Il periodo più lungo della preistoria umana, esteso da circa 3,3 milioni di anni fa fino alla fine dell'ultima glaciazione. È convenzionalmente suddiviso in Paleolitico inferiore (industrie ad amigdale e chopper), Paleolitico medio (industria musteriana, associata in Europa ai Neanderthal) e Paleolitico superiore (industrie lamellari, associata in Europa a \textit{Homo sapiens} a partire da circa 43.000 anni prima dell'anno zero). Nell'Italia nordorientale il Paleolitico medio è attestato nei siti della Lessinia fino a circa 38.000 anni prima dell'anno zero; il Paleolitico superiore, con la cultura epigravettiana, è presente nei siti prealpini veneti fino a circa 11.000 anni prima dell'anno zero.

\item[Paraglaciale] Aggettivo che indica i processi geomorfici instabili e accelerati che seguono il ritiro di un ghiacciaio: frane, colate detritiche, soliflusso, riassestamento delle pareti rocciose private del sostegno del ghiaccio. La causa principale è la perdita del supporto laterale e basale fornito dal ghiaccio, che lascia versanti sovra-ripiditi e materiali sciolti in condizioni di instabilità. La fase paraglaciale ha una durata variabile da decenni a millenni e si conclude quando il paesaggio raggiunge un nuovo equilibrio. Nelle Alpi, la fase paraglaciale principale si sviluppò tra circa 17.000 e 13.000 anni prima dell'anno zero, durante il ritiro dei ghiacciai tardoglaciali.

\item[Qhapaq Ñan] Rete stradale dell'impero inca (\textit{Tawantinsuyu}), che collegava tra loro le principali regioni dell'impero lungo la dorsale andina e i versanti verso il Pacifico e l'Amazzonia, per una lunghezza complessiva di oltre 30.000 chilometri. La costruzione, iniziata da stati predecessori e ampliata sistematicamente dagli Inca a partire dal XV secolo d.C., era sostenuta dal sistema di lavoro collettivo obbligatorio (\textit{mit'a}). La rete comprendeva ponti, terrapieni, gradinate e stazioni di sosta (\textit{tambo}) a intervalli regolari. Dal 2014 è riconosciuta dall'UNESCO come patrimonio mondiale dell'umanità, nella forma di un sito seriale che comprende sezioni in sei paesi sudamericani.

\item[Research bias] Distorsione sistematica che si produce quando la distribuzione geografica dei ritrovamenti scientifici rispecchia la distribuzione degli sforzi di ricerca piuttosto che la distribuzione reale del fenomeno studiato. In archeologia preistorica, si manifesta quando le aree con una tradizione di scavi sistematici mostrano densità di siti molto superiori ad aree comparabili ma meno indagate, non perché il territorio fosse effettivamente più frequentato, ma perché nessuno ha ancora cercato in modo sistematico. Il termine è usato anche in ecologia (distribuzione delle specie note), epidemiologia (incidenza di malattie diagnosticate) e altri campi in cui la conoscenza dipende dall'intensità dell'indagine.

\item[Sauveterriano] Cultura litica che contraddistingue la fase più antica del Mesolitico nell'Italia settentrionale e in buona parte dell'Europa occidentale, collocata approssimativamente tra 12.700 e 9.500 anni prima dell'anno zero. Prende il nome dal sito di Sauveterre-la-Lémance (Lot-et-Garonne, Francia). L'industria è caratterizzata da microliti geometrici di piccole dimensioni (segmenti, triangoli, punte a dorso curvo) ottenuti da supporti lamellari molto regolari. I siti sauveterriani alpini mostrano preferenza per posizioni di visibilità e controllo del territorio, in connessione con la caccia a stambecco e camoscio.

\item[Tambo] Stazione di sosta, rifornimento e controllo amministrativo disposta lungo il Qhapaq Ñan a intervalli di una giornata di cammino, generalmente ogni 20--30 chilometri. I tambo erano strutture permanenti in pietra o adobe, gestite dallo Stato inca, che fornivano alloggio ai funzionari, ai messaggeri (\textit{chasqui}) e ai contingenti militari in transito. Custodivano riserve di cibo, tessuti e altri beni ridistribuiti secondo il sistema della reciprocità asimmetrica inca. L'ubicazione dei tambo non rispondeva esclusivamente a criteri di minimizzazione del percorso, ma anche a necessità amministrative, rituali e strategiche.

\item[Tardoglaciale] Periodo di transizione tra l'Ultimo Massimo Glaciale e l'Olocene, compreso, per le Alpi, tra circa 17.000 e 9.700 anni prima dell'anno zero. È caratterizzato da un andamento climatico oscillante: una prima fase di riscaldamento (con l'oscillazione di Ragogna/Gschnitz tra 15.000 e 13.600 anni prima dell'anno zero), seguita dal miglioramento rapido del Bølling-Allerød (12.700--11.000 anni prima dell'anno zero), interrotto dal raffreddamento del Dryas Recente (11.000--9.700 anni prima dell'anno zero). Durante il Tardoglaciale il ritiro dei ghiacciai alpini rese progressivamente accessibili le quote medie e alte, permettendo la prima frequentazione stabile della fascia montana da parte di cacciatori-raccoglitori.

\item[Ultimo Massimo Glaciale (UMG)] Fase di massima estensione dei ghiacciai durante l'ultima glaciazione, collocata globalmente tra circa 26.000 e 19.000 anni prima dell'anno zero. Nelle Alpi il massimo fu raggiunto intorno a 21.000--19.000 anni prima dell'anno zero: i ghiacciai scendevano fino alle pianure pedemontane, occupando le principali valli alpine fino alle quote di uscita in pianura. Le temperature medie estive erano inferiori di 8--10°C rispetto a oggi; la linea degli equilibri glaciali si trovava circa 1.000--1.200 metri più in basso rispetto all'attuale. La fine dell'UMG, circa 19.000 anni prima dell'anno zero, segna l'inizio del ritiro glaciale e del Tardoglaciale.

\end{description}
